\section{Methodology}
\label{sec:methodology}

We benchmark twelve neural vocoders under one common data, evaluation, and
measurement protocol.  A \emph{Project-Trained Configuration} is one architecture recipe
fitted from random initialization under the project budget; its selected
state is the \emph{Retained Project Checkpoint}, which every measurement
inherits.  The protocol fixes corpus membership, evaluator, and measurement
procedure while each configuration keeps its native conditioning, objective,
and optimization recipe: a structured common testbed, not an equal-compute
causal benchmark.  Throughout, a deployment artifact $a=(m,t,h)$ applies
transformation $t$ to the checkpoint $\theta_m^\star$ of configuration
$m\in\mathcal M$ on requested \emph{hardware lane}
$h\in\{\mathrm{B200},\mathrm{CPU}\}$; the
identity transformation defines the paired same-lane control $c(a)$, and
inference seeds $s\in\{42,43,44\}$ re-execute one selected state without
resampling training (single exception: static ONNX INT8,
Section~\ref{sec:deployment-methods}).

\subsection{Data Set and Preprocessing}
\label{sec:data-preprocessing}

All experiments use LJSpeech~1.1, 13,100 English utterances from one speaker
\citep{ito2017ljspeech}.  Identifier-sorted disjoint blocks fix 12,475
training, 100 validation, and 525 evaluation items without a random seed;
identifier order preserves chapter structure, so the split is fixed
membership rather than a random utterance sample.
Training loaders shuffle and crop fixed-length segments,
whereas validation and evaluation consume complete utterances.  The
evaluation suffix is \emph{adaptive}, not untouched: its PESQ values informed
stop and extension decisions, and one deterministic repair of RFWave's
inference sampler.

Every configuration retains its native conditioning interface: most consume
80 Slaney-scale mel bands at 22.05~kHz; BigVGAN-base uses 100 normalized
bands, and Vocos, VocosFormer, and RFWave use 100 unnormalized HTK-scale
bands at 24~kHz; all mel pipelines share a 1,024-point Hann transform with
a 256-sample hop, while LPCNet consumes 18 Bark-frequency cepstra plus pitch and
correlation features at 16~kHz.  Training-only transformations are peak
normalization to 0.95 (HiFi-GAN V1, V2, and V3), uniform random peak gain in
$[-6,-1]$~dB (Vocos, VocosFormer), and LPCNet excitation noise with standard
deviation 0.3; evaluation waveforms receive none.

\subsection{Model Architectures}
\label{sec:model-architectures}

The cohort spans five generation mechanisms: time-domain adversarial
upsamplers (MelGAN, HiFi-GAN V1, V2, V3, BigVGAN-base), adversarially
trained spectral models with inverse-STFT heads (APNet2, FreeV, RNDVoC,
Vocos), hybrid attention (VocosFormer), sample-level autoregression
(LPCNet), and rectified-flow synthesis (RFWave);
Table~\ref{tab:benchmark-cohort} reports the executed configuration surface,
with the cited architectural definitions in the literature review.

\begin{table}[t]
  \centering
  \footnotesize
  \renewcommand{\arraystretch}{1.02}
  \setlength{\tabcolsep}{4pt}
  \InputIfFileExists{sections/methodology/scripts/generated/cohort_table.tex}{}{%
    \begin{tabular*}{\linewidth}{@{\extracolsep{\fill}}lllrrrcrr@{}}
  \toprule
  Configuration & Family & Conditioning & Params & Segment & Batch & Prec. & Updates & State \\
  \midrule
  APNet2 & iSTFT GAN & 80 mel / 22.05 & 31.43 & 8,192 & 32 & BF16 & 124.8 & 120k \\
  BigVGAN-base & Waveform GAN & 100 mel / 24 & 14.03 & 8,192 & 32 & BF16 & 124.8 & 120k \\
  FreeV & iSTFT GAN & 80 mel / 22.05 & 18.22 & 8,192 & 32 & BF16 & 124.8 & 120k \\
  HiFi-GAN V1 & Waveform GAN & 80 mel / 22.05 & 13.94 & 8,192 & 32 & BF16 & 254.8 & 245k \\
  HiFi-GAN V2 & Waveform GAN & 80 mel / 22.05 & 0.93 & 8,192 & 32 & BF16 & 124.8 & 120k \\
  HiFi-GAN V3 & Waveform GAN & 80 mel / 22.05 & 1.46 & 8,192 & 32 & BF16 & 254.8 & 245k \\
  LPCNet & Autoregressive & 20 feat / 16 & 1.23 & 2,400 & 128 & FP32 & 31.3 & 30k \\
  MelGAN & Waveform GAN & 80 mel / 22.05 & 4.27 & 16,000 & 32 & BF16 & 124.8 & 120k \\
  RNDVoC & iSTFT GAN & 80 mel / 22.05 & 3.57 & 16,384 & 32 & BF16 & 124.8 & 120k \\
  RFWave & Rectified flow & 100 mel / 24 & 18.14 & 32,512 & 64 & BF16 & 125.0 & 120k \\
  Vocos & iSTFT GAN & 100 mel / 24 & 13.53 & 16,384 & 32 & BF16 & 124.8 & 120k \\
  VocosFormer & iSTFT GAN + attn & 100 mel / 24 & 13.78 & 16,384 & 32 & BF16 & 124.8 & 120k \\
  \bottomrule
\end{tabular*}
}
  \caption{Project-Trained Configurations.  Conditioning: input features and
  rate (kHz); Params: deployable parameters (millions); Segment: training
  crop (samples); Updates: executed budget (thousands); State: the Retained
  Project Checkpoint, the terminal durable state selected by the
  registered gate.}
  \label{tab:benchmark-cohort}
\end{table}

VocosFormer, the project-defined exception, inserts a residual-convolution,
self-attention, residual-convolution block from the WavTokenizer decoder
lineage into a shallower Vocos chassis under matched losses, data, and
training quantum; its 13.78 million parameters sit $1.83\%$ above Vocos,
and one training trajectory supports only configuration-level comparison.

\subsection{Training and Checkpoint Selection}
\label{sec:training-checkpoint}

We fitted each configuration through one recorded-seed (1234) training
trajectory under its native objective and a finite registered budget.  All
models are PyTorch modules executed through the project-controlled
SyntheticMind harness, which standardizes fitting, validation, metric
reduction, and checkpoint restore; the VOCODE library
supplies data pipelines, objectives, synthesis paths, evaluators, and
deployment transformations, and each configuration, variant, lane, and
execution runs as one isolated job.
Eleven runs used a 48-GB NVIDIA L40S and RFWave an 80-GB NVIDIA A100;
recorded training sums to 75.31 GPU-hours, a lower bound.
Architecture-native rather than equalized budgets make the retained states
controlled project realizations, not attainable quality ceilings.

Objectives are architecture-native.  MelGAN optimizes least-squares
adversarial and feature-matching terms alone; HiFi-GAN V1, V2, and V3 and
BigVGAN-base combine the same terms over their discriminator ensembles
with weighted mel reconstruction, and Vocos and VocosFormer apply the
hinge form of that composition.  APNet2, FreeV, and
RNDVoC add log-amplitude, anti-wrapping phase, and STFT-consistency terms
with real and imaginary reconstruction to mel, hinge adversarial, and
feature-matching supervision.  LPCNet minimizes teacher-forced cross
entropy over its quantized mu-law excitation; RFWave regresses conditional
flow velocities in the waveform domain with auxiliary magnitude and
band-overlap terms.  Objectives were deliberately not standardized, so raw
training and validation losses are not comparable across families and
never ranked configurations.  Ten configurations optimized with AdamW and
two (MelGAN, LPCNet) with Adam; initial learning rates were
$2{\times}10^{-4}$ (APNet2, FreeV, HiFi-GAN V1, V2, and V3, RNDVoC,
RFWave), $1{\times}10^{-4}$ (BigVGAN-base, MelGAN), $5{\times}10^{-4}$
(Vocos, VocosFormer), and $1{\times}10^{-3}$ (LPCNet).  LPCNet additionally
trains under scheduled GRU sparsification from 2k to 20k updates, the
HiFi-GAN V1 and V3 budgets include 4,800 replayed updates after a durable
resume, and exact coefficients and schedules remain recoverable from the
resolved run records and code.

Checkpoint selection follows a registered gate: validation loss informed
continuation, durable checkpoints were written every 5,000 updates, a
sub-0.02 PESQ gain across gates indicated a plateau, and every run kept a
registered ceiling.  Where a commensurate published or released-checkpoint
PESQ anchor existed, the registered quality target was 93\% of that
anchor; RFWave used a registered PESQ floor of 3.55, and LPCNet and
VocosFormer had no commensurate numeric anchor.  For each configuration, evaluation
used the terminal durable checkpoint admitted by this gate, the Retained
Project Checkpoint $\theta_m^\star$ of Table~\ref{tab:benchmark-cohort},
not a retrospectively substituted validation minimum; the
validation-best, global-minimum, and evaluated states remain in the
resolved records.  The accepted runs predate the repair that seeds model
construction and workers deterministically, so 1234 is a recorded rather
than reconstructible seed; retained states are immutable, and bitwise
regeneration is not claimed.

\subsection{Deployment Optimization Methods}
\label{sec:deployment-methods}

We transform each Retained Project Checkpoint with the deployment variants
of Figure~\ref{fig:deployment-admission}, spanning TorchInductor
compilation, FP16 weight casting, INT8 quantization, ONNX deployment,
magnitude pruning at 30/50/70\% with one recovery arm, and RFWave
solver-step reduction.  Admission requires an executable transformed path,
nonzero operator or parameter coverage, all 525 utterances under three
executions, and a same-profile pre-transformation measurement; unsupported
cells fail before execution expansion.  The executed surface contains 22
pre-transformation controls and 69 transformed groups (91 groups, 273
executions); LPCNet completed no admissible execution on either lane,
leaving eleven transformed configurations, and 29 exclusion records
document 88 rejected attempts.

\begin{figure}[t]
  \centering
  \includegraphics[width=0.88\linewidth]{deployment_admission_matrix}
  \caption{Deployment admission surface.  Filled: admitted three-execution
  groups (91, incl.\ 22 same-lane controls); crosses: the ten excluded
  attempts inside the plotted vocabulary (invalid recovery or control
  states; incomplete LPCNet executions), with the 29-row exclusion register
  remaining authoritative.  Compile RO = reduce-overhead; INT8 WO =
  weight-only; Recovery = 13-epoch recovery of the 50\%-pruned state; ODE
  $k$ = $k$ Euler steps.}
  \label{fig:deployment-admission}
\end{figure}

Static ONNX INT8 uses ONNX
Runtime~1.24.4 on its CPU provider with eight intra-op threads, QDQ format,
per-channel QInt8 weights, QUInt8 activations, and coverage of Conv,
ConvTranspose, MatMul, and Gemm.  Each execution reconstructs its quantized
graph, calibrating on 48 mel batches from the execution-seeded shuffled
training loader, so these executions combine calibration-set and runtime
variation rather than repeating one fixed artifact.  Dynamic and weight-only
INT8 retain their distinct PyTorch and torchao semantics; an unsupported
INT4 probe contributes no result.  Pruning is magnitude-based with dense
masks, testing weight removal but not sparse storage or kernels: Vocos,
FreeV, VocosFormer, and RFWave use $L_2$-ranked structured output-channel
pruning over backbone Linear projections (spectral heads excluded); the rest
use global $L_1$ unstructured pruning over convolutional, Linear, and
recurrent weights.  Selected 50\% states received 13 recovery epochs at one
tenth learning rate, masks held during fine-tuning and zeros baked into the
evaluated checkpoint; absent an admitted dense-continuation control,
recovered rows support descriptive rather than causal claims.  RFWave step
reduction shrinks its ten-step Euler baseline to eight, four, or two steps.

\subsection{Evaluation Metrics and Analysis Protocols}
\label{sec:evaluation-protocols}

\paragraph{Quality proxies.}
Wideband PESQ, withdrawn by ITU-T in January 2024 in favor of the P.863
family \citep{itut2024p862} yet still the dominant quality proxy in
the vocoder literature, remains the primary intrusive measure
\citep{rix2001pesq};
STOI \citep{taal2011stoi} and UTMOS \citep{saeki2022utmos} test dependence on
the proxy choice.  Reference and synthesized waveforms are truncated to
their common unpadded length (\emph{true-length} evaluation) before
resampling; PESQ operates wideband at
16~kHz, classic non-extended STOI at 10~kHz, and UTMOS (SpeechMOS UTMOS22
strong~1.2.0) scores the synthesized waveform alone at 16~kHz.  Mel $L_1$, multi-resolution STFT, mel-cepstral,
log-spectrum, pitch, periodicity, and voicing errors are diagnostics only;
architecture-native or alignment-sensitive values are not cross-model
endpoints, and without a blinded listening study every quality claim remains
a proxy claim.  For utterance $i$ in the fixed 525-item evaluation set
$\mathcal I$, with conditioning $x_i$ and reference $y_i$, artifact $a$
under execution $s$ synthesizes $\widehat y_{a,s,i}=f_{a,s}(x_i)$, and
proxy $q$ is aggregated as
\begin{equation}
  \widehat Q_a^{(q)}
  =\frac{1}{|\mathcal I|}\sum_{i\in\mathcal I}
   \frac{1}{|\mathcal S|}\sum_{s\in\mathcal S}
   q\bigl(y_i,\widehat y_{a,s,i}\bigr),
  \qquad \mathcal S=\{42,43,44\},
  \label{eq:quality-estimator}
\end{equation}
where the reference-free UTMOS ignores $y_i$; execution values are averaged
within each utterance before any paired contrast, and for static ONNX INT8
the execution index also carries calibration-state variation
(Section~\ref{sec:deployment-methods}).

\paragraph{Timing and footprint protocol.}
Beyond quality, endpoints include the continuous process-RSS high-water
mark, one-time ONNX session-construction time, and complexity as
parameters, serialized bytes, and MACs from one traced second of audio;
jobs were independently
scheduled, and paired contrasts match the requested resource profile,
never the physical host.  After five warm-up batches,
pre-transformation evaluation measures one complete pass at batch one on
an NVIDIA B200.  The CPU column of Table~\ref{tab:baseline-surface} is the
retained comparative record, three executions each timing one complete
post-warm-up pass on an 8-core CPU host, and is never cross-compared with
deployment-phase RTFs; LPCNet's 27-item timing record is excluded from
strict speed frontiers.  Deployment evaluation measures three synchronized
calls per utterance for the 520 post-warm-up utterances on B200 or CPU,
each timed call
following one untimed prediction of the same input, so the quantities
describe warm repeated synthesis only.  For artifact $a$, execution $s$,
timed items $\mathcal I_a\subseteq\mathcal I$, synchronized time
$\tau_{a,s,i,r}$, duration $d_i$, and $J_a\in\{1,3\}$ calls, warm RTF is
\begin{equation}
  \bar\tau_{a,s,i}=J_a^{-1}\sum_{r=1}^{J_a}\tau_{a,s,i,r},\qquad
  R_{a,s}=\frac{\sum_{i\in\mathcal I_a}\bar\tau_{a,s,i}}
                 {\sum_{i\in\mathcal I_a}d_i},
  \label{eq:rtf}
\end{equation}
and the group-level $\widehat R_a$ is the mean of $R_{a,s}$ over
executions.  Reported p50/p95 are per-execution percentiles of the 520
per-utterance three-repetition means $\bar\tau_{a,s,i}$, averaged over
executions, not of the 1,560 individual calls, whose raw durations were
not retained.  B200 jobs requested 16 CPUs and 64~GiB, CPU
jobs 8 CPUs and 16~GiB; the Python~3.14.2, PyTorch~2.13.0+cu130,
torchaudio~2.11.0, torchao~0.17.0, ONNX~1.20.1, and ONNX Runtime~1.24.4 stack
is retained, whereas host-identifying fields are irrecoverable.

\paragraph{Paired contrasts and decision rule.}
For quality proxy $q$ and transformed artifact $a$ with paired control
$c(a)$, quality estimates $\widehat Q$ from
Eq.~\eqref{eq:quality-estimator}, warm RTF $\widehat R$ from
Eq.~\eqref{eq:rtf}, and serialized deployable size $D$, the continuous
effects are
\begin{equation}
  \Delta Q_a^{(q)}=\widehat Q_a^{(q)}-\widehat Q_{c(a)}^{(q)},\qquad
  S_a=\widehat R_{c(a)}/\widehat R_a,\qquad
  Z_a=D_a/D_{c(a)}.
  \label{eq:deployment-effects}
\end{equation}
Quality deltas are utterance-paired, whereas speedups and size ratios are
same-profile descriptive contrasts because hosts are not paired.  Quality
contrasts pair the execution-averaged utterance values of
Eq.~\eqref{eq:quality-estimator} before a Student $t$ interval, with a
5,000-replicate percentile bootstrap as sensitivity analysis
\citep{student1908probableerror,efron1979bootstrap}.  These are fixed-set
descriptive intervals quantifying utterance dispersion over the 525 fixed
items under a working exchangeability assumption; they omit multiplicity
correction and variation from training, speaker, recovery seed, or host,
and for static INT8 do not separate calibration from runtime variation.
A post hoc ordered classifier provides an exploratory engineering summary
of Eq.~\eqref{eq:deployment-effects}: dense-masked pruning and recovery rows
are labeled first: harmful when $\Delta Q_a^{(\mathrm{PESQ})}<-0.10$,
otherwise inconclusive; remaining MelGAN rows are fixed as inconclusive
because their PESQ is floor-censored.  Every other row receives
\begin{equation}
  \ell(a)=
  \begin{cases}
    \text{harmful}, &
      \Delta Q_a^{(\mathrm{PESQ})}<-0.10
      \ \lor\ \bigl(S_a<0.90\wedge Z_a>0.90\bigr)\\[1pt]
    \text{beneficial}, &
      \Delta Q_a^{(\mathrm{PESQ})}\ge-0.05
      \ \wedge\ S_a\ge0.90
      \ \wedge\ \bigl(S_a\ge1.10\lor Z_a\le0.55\bigr)\\[1pt]
    \text{inconclusive}, & \text{otherwise.}
  \end{cases}
  \label{eq:decision-rule}
\end{equation}
Thresholding PESQ alone makes these labels PESQ-preserving statements, not
perceptual equivalence claims.  A $3^5=243$ grid tests label sensitivity by
varying each cutoff in Eq.~\eqref{eq:decision-rule} over three levels
(preservation $-0.075/-0.05/-0.025$; harm $-0.15/-0.10/-0.075$; speed floor
$0.85/0.90/0.95$; gain $1.05/1.10/1.20$; size $0.50/0.55/0.60$).

\paragraph{Frontier and constraint-aware selection.}
Under one quality proxy $q$ and one cost axis (RTF, parameter count, or
serialized size), artifact $a$ dominates $a'$ when
$\widehat Q_a^{(q)}\ge\widehat Q_{a'}^{(q)}$ at no greater cost, with at
least one inequality strict; each reported frontier is the nondominated
set, recomputed per proxy and axis.  For
hardware lane $h$, admitted artifacts $\mathcal A_h$, quality floor
$q_{\min}$, parameter ceiling $P_{\max}$, and size ceiling $D_{\max}$,
deployment selection solves
\begin{equation}
  a_h^\star\in\underset{a\in\mathcal A_h}{\arg\min}\ \widehat R_a
  \quad\text{s.t.}\quad
  \widehat Q_a^{(q)}\ge q_{\min},\quad P_a\le P_{\max},\quad
  D_a\le D_{\max},
  \label{eq:constrained-selection}
\end{equation}
so every recommendation is conditional on hardware, proxy, and constraints.
